\subsection{An independent proof through degree sixteen, at every shift}
There is another route for this particular Artin factor that uses the
location of its low zeros. Define $G(w)=\Lambda_f(1/2+\sqrt w)$, with
analyticity at zero given by evenness. We prove
\[
 J^{d,n}(X)=\sum_{j=0}^d \binom{d}{j}\gamma_{n+j}X^j
 \quad\text{has only negative real zeros for } 1\le d\le16,\quad n\ge0 .
\]
There are two ingredients.

\textit{A zero-free rectangle.}
Put $a=\Re(s)-1/2$, $t=\Im(s)$. For the conductor-$23$ eta kernel of Section~\ref{sec:eta-jensen},
\[
 \tfrac12\Re\Lambda_f(s)=\int_0^\infty K(u)\cosh(au)\cos(tu)\,du.
\]
On $|a|\le1/2$ and $|t|\le4$, the accompanying \texttt{check\_height.py}
proves the right side exceeds $1/1000$ using integer bounds.
Here are the steps of the certificate. Subdivide $0\le u\le4$ in steps
of $1/1024$ and $0\le t\le4$ in steps of $1/256$. In each $u$ interval
bound $q=\exp[-(2\pi/\sqrt{23})e^u]$ using
$1.3101<2\pi/\sqrt{23}<1.3102$. In
$K=e^{u/2}q\prod_{m\ge1}(1-q^m)(1-q^{23m})$ keep the first eight
factors of the first product. The product of the remaining factors is
between $1-2q^9/(1-q)$ and $1$, by
$\prod(1-v_i)\ge1-\sum v_i$ for $0\le v_i\le1$. Bound the phase cosine
from below by evaluating its degree-$62$ lower Taylor polynomial at
$\lfloor512 t_{\min}u_{\min}\rfloor/512$ and subtracting the maximum
phase distance from this point in the cell. This uses $|\cos'|\le1$.
A positive cosine lower bound is multiplied by the lower bound for
$K$, a negative one by the upper bound for $K\cosh(u/2)$.
The sum of these interval contributions is a lower integral for every $a$
and every $t$ in that cell of heights. Beyond $u=4$ the absolute integrand
is bounded by $e^u e^{-e^u}$, giving tail $e^{-e^4}<2^{-64}$.
All exponentials in the certificate use a rational Taylor bound on
$[0,1]$ followed by seven squarings with outward integer rounding.
Thus no floating-point or library zero-finder is required.

\textit{The sector argument.}
The Euler product and functional equation put the zeros of the completed
Artin function in $|a|\le1/2$. The zeros of $G$ are $w=(a+it)^2$.
If all zeros with $|t|<T$ are on the line $a=0$, or merely satisfy
$|a|/|t|\le1/(2T)$, the zeros of $G$ lie in the closed sector about
the negative real ray with half-angle $\theta=2\arctan(1/(2T))$,
provided $T>1/2$ and $\Lambda_f(1/2)\ne0$. The growth estimate
$\log\max_{|z|\le R}|\Lambda_f(1/2+z)|=O(R\log(2+R))$ follows by
integrating $e^{R|u|-c e^{|u|}}$. Hence $G$ has order at most $1/2$ and
is the locally uniform limit of its genus-zero partial products. Group
conjugate zeros to obtain real polynomials. The closed sector is convex;
Gauss--Lucas and locally uniform convergence place the zeros of every
$G^{(n)}$ in the same sector. All $G^{(n)}$ have positive Taylor
coefficients by the Mellin formula.

The sector theorem of Obreschkoff and Chasse now gives real-rooted
Jensen polynomials for each $G^{(n)}$ through degree
\[
   d\le\csc^2\theta=\left(T+\frac{1}{4T}\right)^2.
\]
See M. Chasse, \emph{Laguerre multiplier sequences and sector properties
of entire functions}, Complex Var. Elliptic Equ. 58 (2013), 875--885,
\href{https://doi.org/10.1080/17476933.2011.584250}{journal version}, or
Theorem~171 and the application to derivatives in Theorem~173 of his
\href{https://people.kth.se/~chasse/dissertation.pdf}{dissertation}.
Taking $T=4$ gives $d\le\lfloor(65/16)^2\rfloor=16$.
Positivity of the coefficients places all real roots on the negative ray.

\textit{Larger ranges from external computations.}
The same sector calculation with $T=18.08$ gives $d\le327$ from the
published interval data. LMFDB identifies this function with classical
newform $23.1.b.a$, lists ten positive low ordinates from $5.1156833$
through $18.0806617$, and records rank zero. Its documentation for modular
forms describes a complete list of zero disks within the listed range.
The displayed radii give $|a|<10^{-3}$ and $|t|>5$ in those disks.
This suffices for $|a|/|t|<1/(2T)$ even without a proof that each disk's
zero lies exactly on the line. See the
\href{https://www.lmfdb.org/L/2/23/23.22/c0/0/0}{record} and
\href{https://www.lmfdb.org/L/2/23/23.22/c0/0/0/Reliability}{precision
and completeness statement}.

In his thesis Nathan Ng reports $1{,}736$ positive zeros and verification
of RH for this Artin factor up to $T=972.2$ using Tollis's zero-counting
theorem. Taking that computation as input gives $d\le945{,}173$. The
thesis illustrates the Gram-block test in detail at level $31$ and
tabulates its outcome for level $23$. We did not rerun the height-$972.2$
computation. See \emph{Limiting Distributions and Zeros of Artin
$L$-Functions}, \href{https://www.cs.uleth.ca/~nathanng/RESEARCH/phd.thesis.pdf}{Secs.~4.3.4--4.3.5}.
None of these finite ranges uses the cubic kernel assertion above.
